\section{第 1 讲：随机对照试验}
\sourcepages{Lecture1-10.pdf，第 2、1、3 页；阅读顺序：先第 2 页正文，再第 1 页补充，最后第 3 页。}
在随机对照试验（RCT）中，用回归分析平均处理效应，是为了减小方差；ATE 的识别并不依赖回归。
\notearrow{black}{RCT 中的回归（第 2 页）}{减小方差；ATE 的识别来自随机分配。}

\[
\Delta_i=Y_i(1)-Y_i(0),\qquad \tau=E[Y_i(1)-Y_i(0)].
\]
用均值差估计量估计 $\tau$，并进行统计推断：
\[
\widehat\tau_{\rm DM}=\frac1{n_1}\sum_{W_i=1}Y_i-\frac1{n_0}\sum_{W_i=0}Y_i,
\qquad n_w=|\{i:W_i=w\}|.
\]
笔记标注：该估计量是一致且无偏的。
\notearrow{blue}{均值差（第 2 页）}{$\widehat\tau_{\rm DM}$}

\begin{reviewnote}{qualification}{S1-002}{抽样与处理分配条件}{Lecture1-10.pdf，第 2 页}
在上述随机分配设定下，两组均非空时，均值差是无偏的；总体一致性与正态极限还需要独立同分布抽样及有限二阶矩。下面的平衡分配方差公式专指 $P(W=1)=1/2$，并假设协变量已中心化，即 $E[X]=0$。一般方差为 $\operatorname{Var}(Y(1))/e+\operatorname{Var}(Y(0))/(1-e)$。中心化也解释了后面的 OLS 分解。在仅有上述矩条件时，分解旁的 $N(0,\cdot/n_w)$ 是渐近近似，而非有限样本下的精确分布。参见讲义式 (1.6)、(1.8)--(1.13)。
\end{reviewnote}

\paragraph{补充：假设。}
\[
Y_i=Y_i(W_i),\qquad W_i\perp\{Y_i(0),Y_i(1)\}.
\]
第一个式子标为 SUTVA，第二个是处理的随机分配。SUTVA 即稳定单位处理值假设；原笔记简述为：（1）单位稳定，实验单位之间相互独立；（2）只有一个处理水平／版本。
\notearrow{black}{$Y_i=Y_i(W_i)$；SUTVA}{稳定单位处理值假设：左页补充中的两项解释（第 1 页）。}

\begin{reviewnote}{original}{S1-001}{SUTVA 不等于统计独立}{Lecture1-10.pdf，第 1 页}
SUTVA 要求单位之间没有干扰，且每种指定处理没有隐藏版本。它不要求结果在统计上相互独立，也允许这里的 $0,1$ 两种处理取值。独立抽样是另外的假设。原笔记的“单一处理水平”宜改为“每一种处理都有定义明确的版本”。参见\href{https://web.stanford.edu/~swager/stats361.pdf}{讲义}印刷页码第 2--3 页的讨论与脚注。
\end{reviewnote}

\subsection{线性结果模型与方差缩减}
\[
Y_i(w)=c_{(w)}+X_i\beta_{(w)}+\varepsilon_i(w),\quad
E[\varepsilon_i(w)\mid X_i]=0,\quad
\operatorname{Var}[\varepsilon_i(w)\mid X_i]=b^2.
\]
\inkblue{两种处理状态分别拟合各自的系数。}记 $A=\operatorname{Var}[X]$，\inkblue{$\|v\|_A^2=v'Av$}。笔记中的方差为
\[
V_{\rm DM}=4b^2+\|\beta_{(0)}+\beta_{(1)}\|_A^2
 +\|\beta_{(0)}-\beta_{(1)}\|_A^2.
\]
另有
\[
\tau=c_{(1)}-c_{(0)}+E[X](\beta_{(1)}-\beta_{(0)}),\qquad
\widehat\tau_{\rm OLS}=\widehat c_{(1)}-\widehat c_{(0)}
 +\overline X(\widehat\beta_{(1)}-\widehat\beta_{(0)}).
\]
笔记给出的误差分解为
\begin{align*}
\widehat\tau_{\rm OLS}-\tau={}&(\widehat c_{(1)}-c_{(1)})-(\widehat c_{(0)}-c_{(0)})
 +\overline X(\beta_{(1)}-\beta_{(0)})\\
&+\overline X(\widehat\beta_{(1)}-\widehat\beta_{(0)}-\beta_{(1)}+\beta_{(0)}).
\end{align*}
前两项分别标注为 $N(0,b^2/n_1)$、$N(0,b^2/n_0)$；第三项标注为 $N(0,\|\beta_{(1)}-\beta_{(0)}\|_A^2/n)$；最后一项是 $O_p(1/n)$。\inkblue{页边批注：斜率估计带来的不确定性趋于零。}因此
\[
\sqrt n(\widehat\tau_{\rm OLS}-\tau)\Rightarrow N(0,V_{\rm OLS}),\qquad
V_{\rm OLS}=4b^2+\|\beta_{(0)}-\beta_{(1)}\|_A^2,
\]
\[
V_{\rm DM}=V_{\rm OLS}+\|\beta_{(0)}+\beta_{(1)}\|_A^2.
\]

\subsection{一般结果回归下的 OLS}
下一页说明：在该模型中，OLS 比均值差更好。设
\[
Y_i(w)=\mu_{(w)}(X_i)+\varepsilon_i(w),\quad
E[\varepsilon_i(w)\mid X_i]=0,\quad
\operatorname{Var}[\varepsilon_i(w)\mid X_i]=b^2.
\]
对于均值差估计量，
\[
\sqrt n(\widehat\tau_{\rm DM}-\tau)\Rightarrow N(0,V_{\rm DM}),\quad
V_{\rm DM}=4b^2+2\operatorname{Var}[\mu_{(0)}(X_i)]
+2\operatorname{Var}[\mu_{(1)}(X_i)].
\]
对于 OLS，Huber--White 方差计算从总体投影出发：
\[
(c^*_{(w)},\beta^*_{(w)})=arg\min_{c,\beta}
 E[(Y_i(w)-X_i\beta-c)^2].
\]
笔记写为
\[
\sqrt{n_w}\left\{\begin{pmatrix}\widehat c_{(w)}\\\widehat\beta_{(w)}\end{pmatrix}
-\begin{pmatrix}c^*_{(w)}\\\beta^*_{(w)}\end{pmatrix}\right\}
\Rightarrow N\left(0,\begin{pmatrix}\operatorname{MSE}^*_{(w)}&0\\0&\cdots\end{pmatrix}\right),
\]
其中 $c^*_{(w)}=E[Y_i(w)]$，并定义
\begin{reviewnote}{original}{S1-003}{错设回归的协方差矩阵不一定为分块对角阵}{Lecture1-10.pdf，第 3 页}
非对角块为零不能由 $E[X]=0$ 推出。令 $q=(1,X')'$、$r_w=Y(w)-c_w^*-X\beta_w^*$、$Q=E[qq']$，则协方差矩阵为
\[
Q^{-1}E[r_w^2qq']Q^{-1}.
\]
截距与斜率之间的协方差块为 $E[r_w^2X]A^{-1}$（行列方向相应调整）；条件均值模型错设时，它一般不为零。投影正交性给出的是 $E[r_wX]=0$，而非 $E[r_w^2X]=0$。截距方差 $E[r_w^2]$ 和随后平衡试验下的方差缩减恒等式仍成立。这是原扫描中可见的错误，讲义式 (1.16) 也有同样问题，并非识别错误。
\end{reviewnote}

\[
\operatorname{MSE}^*_{(w)}=E[(Y_i(w)-X_i\beta^*_{(w)}-c^*_{(w)})^2].
\]
由此得到
\[
\sqrt n(\widehat\tau_{\rm OLS}-\tau)\Rightarrow N(0,V_{\rm OLS}),\qquad
V_{\rm OLS}=V_{\rm DM}-\|\beta^*_{(0)}+\beta^*_{(1)}\|_A^2.
\]
\inkred{$V_{\rm OLS}<V_{\rm DM}$ 总是成立。}
\begin{reviewnote}{original}{S1-004}{方差不会增大，但不一定严格减小}{Lecture1-10.pdf，第 3 页}
前面的恒等式只能推出 $V_{\rm OLS}\le V_{\rm DM}$。当 $\|\beta_0^*+\beta_1^*\|_A^2=0$ 时等号成立；若 $A$ 非奇异，等价于 $\beta_0^*+\beta_1^*=0$。两组斜率非零但互为相反数，就会出现相等的情况。严格缩减要求上述平方范数为正。这可以直接由两个方差相减得到。
\end{reviewnote}


\paragraph{个人理解。}
OLS 可以看作
\[
\widehat\tau_{\rm OLS}=\frac1n\sum_{i=1}^n
\{(\widehat c_{(1)}+X_i\widehat\beta_{(1)})
 -(\widehat c_{(0)}+X_i\widehat\beta_{(0)})\}
=\frac1n\sum_i[\widehat\mu_{(1)}(X_i)-\widehat\mu_{(0)}(X_i)].
\]
\inkblue{这里 $\widehat\mu_{(w)}(x)$ 是 OLS 在 $x$ 处的拟合值，也可以换成其他方法。}

\section{第 2 讲：无混杂性与倾向得分}
\sourcepages{Lecture1-10.pdf，第 5、4、7、6 页；第 4 页承接第 5 页底部，第 6 页补充第 7 页。}
无混杂性：目的是消除混杂因素的影响。倾向得分：给定协变量后接受处理的概率。
\subsection{汇总各组的均值差估计}
\inkred{直接按照概率加权。}对离散且取数值的 $X_i$，假设
\[
\{Y_i(0),Y_i(1)\}\perp W_i\mid X_i=x,\qquad x\in\mathcal X.
\]
\inkblue{这里 $x$ 是分组变量。}定义各组的 ATE：
\[
\tau(x)=E[Y_i(1)-Y_i(0)\mid X_i=x].
\]
则
\begin{align*}
\widehat\tau_{\rm AGG}&=\sum_{x\in\mathcal X}\frac{n_x}{n}\widehat\tau(x),\\
\widehat\tau(x)&=\frac1{n_{x1}}\sum_{X_i=x,W_i=1}Y_i
-\frac1{n_{x0}}\sum_{X_i=x,W_i=0}Y_i,\\
n_x&=|\{i:X_i=x\}|,\qquad
n_{xw}=|\{i:X_i=x,W_i=w\}|.
\end{align*}
令 $e(x)=P(W_i=1\mid X_i=x)$。组内的极限分布为
\begin{align*}
\sqrt{n_x}(\widehat\tau(x)-\tau(x))
&\Rightarrow N\left(0,
\frac{\operatorname{Var}[Y_i(0)\mid X_i=x]}{1-e(x)}
+\frac{\operatorname{Var}[Y_i(1)\mid X_i=x]}{e(x)}\right)\\
&=N\left(0,\frac{b^2(x)}{e(x)(1-e(x))}\right),
\end{align*}
第二个等号使用了原笔记中两种处理状态具有相同条件方差的假设。原定义一行写作 $\widehat\pi(x)=P(X_i=x)=n_x/n$。
\begin{reviewnote}{original}{S1-005}{总体分层概率与经验分层概率}{Lecture1-10.pdf，第 5 页}
总体概率质量应记为 $\pi(x)=P(X_i=x)$，其估计为 $\widehat\pi(x)=n_x/n$；有限样本中二者通常不相等。后面的 $\widehat\pi(x)-\pi(x)$ 已经要求区分二者。以下展开按这两个不同定义理解。固定分层数下的中心极限定理还需要：层数有限、各层总体概率为正、层内处理重叠，以及有限条件二阶矩。参见讲义式 (2.3)--(2.7)。
\end{reviewnote}
于是
\begin{align*}
\widehat\tau_{\rm AGG}={}&\sum_x\pi(x)\tau(x)
+\sum_x\pi(x)(\widehat\tau(x)-\tau(x))\\
&+\sum_x(\widehat\pi(x)-\pi(x))\tau(x)
+\sum_x(\widehat\pi(x)-\pi(x))(\widehat\tau(x)-\tau(x)).
\end{align*}
\inkblue{第一项是 $\tau$。}\inkblue{第二项标注为渐近正态，其方差为}
\begin{inkblock}{blue}
\[
\sum_x\pi^2(x)\operatorname{Var}[\widehat\tau(x)].
\]
\end{inkblock}
\inkblue{第三项的方差是 $n^{-1}\operatorname{Var}[\tau(X_i)]$，最后一项为 $O_p(1/n)$。}沿箭头接到左页，
\notearrow{black}{汇总估计量的四项展开（第 5 页）}{汇总估计量的方差 $V_{\rm AGG}$（第 4 页）。}

\[
\sqrt n(\widehat\tau_{\rm AGG}-\tau)\Rightarrow N(0,V_{\rm AGG}),
\]
\begin{align*}
V_{\rm AGG}&=\operatorname{Var}[\tau(X_i)]
+\sum_{x\in\mathcal X}\pi^2(x)\frac1{\pi(x)}
\frac{b^2(x)}{e(x)(1-e(x))}\\
&=\operatorname{Var}[\tau(X_i)]
+E\left[\frac{b^2(X_i)}{e(X_i)(1-e(X_i))}\right].
\end{align*}
\subsection{连续协变量与倾向得分}
给定 $X_i$ 后，$W_i$ 不能“偷看”$Y_i$：这是作者对无混杂性的直觉理解。若
\[
\{Y_i(0),Y_i(1)\}\perp W_i\mid X_i,
\quad e(x)=P(W_i=1\mid X_i=x),
\]
则
\[
\{Y_i(0),Y_i(1)\}\perp W_i\mid e(X_i).
\]
\paragraph{补充图示与计算。}
\begin{center}
\begin{tikzpicture}[x=1cm,y=.32cm,>=Stealth,font=\small]
\draw[->] (0,0)--(4,0);\draw[->](0,0)--(0,11);
\foreach \x/\u/\l/\d in {1/10/5/5,2/8/4/4,3/7/5/2}{
\fill[blue](\x,\u) circle(2pt);\fill[red](\x,\l) circle(2pt);
\node[above] at (\x,\u){$\u$};\node[below]at(\x,\l){$\l$};
\draw[->](\x,\u-.6)--node[right]{$\d$}(\x,\l+.6);
\node[below]at(\x,0){$x_{\x}$};}
\node at(2,12){离散 $X$};
\begin{scope}[xshift=6cm]
\draw[->](0,0)--(4,0);\draw[->](0,0)--(0,11);
\draw[blue,thick](.7,9).. controls(1.8,10)and(2.5,5.5)..(3.2,6.5);
\draw[red,thick](.7,2).. controls(1.5,1)and(2.5,6)..(3.2,4.5);
\draw[->](.5,8)--(.5,3.5);\node at(2,12){连续 $X$};
\node[below] at(2,0){$x\in\mathcal X\quad\longleftrightarrow\quad e(x)$};
\end{scope}
\end{tikzpicture}
\end{center}
图下方的证明写作
\begin{align*}
&P(W_i=w\mid\{Y_i(0),Y_i(1)\},e(X_i))\\
&\quad=\int_{\mathcal X}P(W_i=w\mid\{Y_i(0),Y_i(1)\},X_i=x)
 P(X_i=x\mid e(X_i))\,dx\\
&\quad=\int_{\mathcal X}P(W_i=w\mid X_i=x)P(X_i=x\mid e(X_i))\,dx\\
&\quad=e(X_i)\mathbf1_{\{w=1\}}+(1-e(X_i))\mathbf1_{\{w=0\}}.
\end{align*}
\inkblue{页边标注：“Lebesgue 积分”。}
\begin{reviewnote}{original}{S1-006}{倾向得分证明中的条件集合}{Lecture1-10.pdf，第 6 页}
令 $\mathcal Y=(Y(0),Y(1))$。迭代期望中的积分一开始必须使用 $dP(X\mid\mathcal Y,e(X))$，而非 $dP(X\mid e(X))$。无混杂性把被积函数化为 $P(W=w\mid X)$；它在每个倾向得分的等值集合上为常数。无论剩余条件分布是什么，对这一常数积分都得到所需结论。对于连续 $X$，应使用正则条件概率测度，而不能把点概率 $P(X=x\mid\cdot)\,dx$ 当作积分测度。可对照讲义命题 2.1 中条件集合正确的证明。
\end{reviewnote}


\subsection{按倾向得分分层}
\begin{enumerate}
\item 排序，使 $\widehat e(X_{i_1})\leq\widehat e(X_{i_2})\leq\cdots\leq\widehat e(X_{i_n})$。
\item 将样本分为 $J$ 层，在每一层内估计均值差。若 $I_j$ 包含排序后第 $\lfloor(j-1)n/J\rfloor+1$ 到第 $\lfloor jn/J\rfloor$ 个样本，则原页的简写公式可表示为
\[
\widehat\tau_j=\frac{\sum_{i\in I_j}W_iY_i}{\sum_{i\in I_j}W_i}
-\frac{\sum_{i\in I_j}(1-W_i)Y_i}{\sum_{i\in I_j}(1-W_i)}.
\]
\notearrow{blue}{使处理组与对照组更相似（第 7 页）}{按 $\widehat e(X_i)$ 排序，并按倾向得分分层。}
\item 对各层估计取平均：$\widehat\tau_{\rm STRAT}=J^{-1}\sum_{j=1}^J\widehat\tau_j$。\inkblue{笔记标注其具有一致性。}
\begin{reviewnote}{qualification}{S1-007}{分层区间需要逐渐缩小}{Lecture1-10.pdf，第 7 页}
若只用固定数量的宽倾向得分区间，一般仍会留下混杂。一致性要求分层越来越细，每层有足够的处理组和对照组样本，满足重叠性，且倾向得分估计足够准确。讲义明确要求 $\widehat e$ 一致收敛，并且 $J$ 随 $n$ 以适当速度增长。这里各层样本量近似相等，因此使用等权平均。
\end{reviewnote}

\end{enumerate}
\subsection{逆倾向得分加权}
\sourcepages{Lecture1-10.pdf，第 9、8、10 页；第 8 页给出第 9 页方差计算所用的模型。}
\[
\widehat\tau_{\rm IPW}=\frac1n\sum_{i=1}^n
\left\{\frac{W_iY_i}{\widehat e(X_i)}-\frac{(1-W_i)Y_i}{1-\widehat e(X_i)}\right\}.
\]
理想的 IPW（oracle IPW）假定完全知道 $e(x)$：
\[
\widehat\tau^*_{\rm IPW}=\frac1n\sum_i
\left\{\frac{W_iY_i}{e(X_i)}-\frac{(1-W_i)Y_i}{1-e(X_i)}\right\}.
\]
假设满足重叠性：所有 $x\in\mathcal X$ 上均有 $\eta\leq e(x)\leq1-\eta$，且 $|Y_i|<M$。若
\[
\sup_{x\in\mathcal X}|e(x)-\widehat e(x)|=O_p(a_n)\longrightarrow0,
\]
笔记给出
\[
|\widehat\tau_{\rm IPW}-\widehat\tau^*_{\rm IPW}|=O_p(a_nM/\eta).
\]
\inkblue{因此，理想估计量的一致性可以传递给使用估计倾向得分的 IPW。}另外，
\[
E[\widehat\tau^*_{\rm IPW}]=E[Y_i(1)-Y_i(0)]=\tau.
\]
\inkblue{满足重叠性时，理想 IPW 的收敛尺度是 $1/\sqrt n$。}
\paragraph{方差计算：左页补充的模型。}
左页给出
\begin{align*}
Y_i(0)&=c(X_i)-(1-e(X_i))\tau(X_i)+\varepsilon_i(0),\\
Y_i(1)&=c(X_i)+e(X_i)\tau(X_i)+\varepsilon_i(1),\\
E[\varepsilon_i(w)\mid X_i]&=0,\qquad
\operatorname{Var}[\varepsilon_i(w)\mid X_i=x]=b^2(x).
\end{align*}
随后笔记写为
\notearrow{black}{第 8 页补充的结果模型}{第 9 页理想 IPW 的方差 $V_{\rm IPW^*}$。}

\begin{align*}
n\operatorname{Var}[\widehat\tau^*_{\rm IPW}]
=V_{\rm IPW^*}
={}&E\left[\frac{c^2(X_i)}{e(X_i)(1-e(X_i))}\right]
+\operatorname{Var}[\tau(X_i)]\\
&+E\left[\frac{b^2(X_i)}{e(X_i)(1-e(X_i))}\right],\\
\sqrt n(\widehat\tau^*_{\rm IPW}-\tau)&\Rightarrow N(0,V_{\rm IPW^*}).
\end{align*}
\paragraph{个人比较。}
$c(x)$ 描述的是 $X$ 与 $Y$ 的关系，而不是 $X$ 与 $\Delta Y$ 的关系。与分组汇总相比，
\[
V_{\rm IPW^*}=V_{\rm AGG}+E\left[\frac{c^2(X_i)}{e(X_i)(1-e(X_i))}\right].
\]
\notearrow{blue}{$V_{\rm IPW^*}-V_{\rm AGG}$（第 10 页）}{\inkblue{除非 $c(X_i)$ 处处为零，否则理想 IPW 的方差比汇总估计更大。}}与线性模型相比：IPW 基于 $\{Y_i(0),Y_i(1)\}\perp W_i\mid X_i$，线性模型则依赖线性关系。
\begin{reviewnote}{qualification}{S1-009}{方差比较与因果识别}{Lecture1-10.pdf，第 10 页}
在重叠性和有限矩条件下，恰当且仅当 $c(X)=0$ 几乎处处成立时，额外方差为零；若 $P(c(X)\ne0)>0$，它就为正。正确设定的线性结果模型是估计上的限制，本身不是因果识别假设：观察性回归仍需要适当的外生性／无混杂性，或其他识别依据。这里保留作者对两种估计方法的比较。
\end{reviewnote}



\section{第 3 讲：通过增广 IPW 高效估计处理效应}
\sourcepages{Lecture1-10.pdf，第 11--12 页。}
普通 IPW 的大样本性质较差。例如，$X$ 离散时，它不如汇总估计量 $\widehat\tau$。
\subsection{ATE 的两种表征}
\begin{enumerate}
\item 利用 $e(x)=P(W_i=1\mid X_i=x)$，有 $\tau=E[\widehat\tau^*_{\rm IPW}]$；理想估计量定义如上。
\item 利用 $\mu_{(w)}(x)=E[Y_i(w)\mid X_i=x]$，
\[
\tau(x)=\mu_{(1)}(x)-\mu_{(0)}(x),\quad
\tau=E[\mu_{(1)}(X_i)-\mu_{(0)}(X_i)],
\]
由此得到
\[
\widehat\tau_{\rm REG}=\frac1n\sum_i[\widehat\mu_{(1)}(X_i)-\widehat\mu_{(0)}(X_i)].
\]
\end{enumerate}
\subsection{增广 IPW 估计量}
\begin{align*}
\widehat\tau_{\rm AIPW}=\frac1n\sum_i\bigg[&\widehat\mu_{(1)}(X_i)-\widehat\mu_{(0)}(X_i)
+W_i\frac{Y_i-\widehat\mu_{(1)}(X_i)}{\widehat e(X_i)}\\
&-(1-W_i)\frac{Y_i-\widehat\mu_{(0)}(X_i)}{1-\widehat e(X_i)}\bigg].
\end{align*}
\paragraph{个人直觉。}
\inkblue{先用 $\mu(x)$ 估计结果，再对回归残差做 IPW。IPW 让数据的协变量构成变得“公平”。}
\subsection{双重稳健性}
只要结果回归 $\widehat\mu_{(w)}(x)$ 或倾向得分 $\widehat e(x)$ 中有一方一致，AIPW 就一致。
\begin{reviewnote}{qualification}{S1-010}{双重稳健性和效率需要正则条件}{Lecture1-10.pdf，第 11、12、13 页}
双重稳健性需要重叠性、适当的可积性，以及对拟合函数求值的控制（例如交叉拟合）；另一个辅助估计量也必须表现良好。仅条件于 $X_i,W_i$，不能把使用同一个样本点训练出来的函数视为固定。交叉拟合的证明还需要条件于训练折。下面的效率界适用于非参数无混杂模型中的正则估计量，而不是字面上所有可能的估计量。下一页列出的乘积风险条件与辅助函数一致性，是选定学习方法后需要验证的假设，并非所有机器学习算法自动具备的性质。参见讲义式 (3.8) 及式 (3.16) 后的正则性讨论。
\end{reviewnote}

当 $\widehat\mu_{(w)}(x)\approx\mu_{(w)}(x)$ 时，回归均值是一致的，加权残差项的均值近似为零，因为
\[
E[Y_i-\widehat\mu_{(W_i)}(X_i)\mid X_i,W_i]\approx0.
\]
这里通过条件化把倾向得分项固定下来。若是 $\widehat e(x)\approx e(x)$，则改写为
\begin{align*}
\widehat\tau_{\rm AIPW}={}&\frac1n\sum_i\left[\frac{W_iY_i}{\widehat e(X_i)}
-\frac{(1-W_i)Y_i}{1-\widehat e(X_i)}\right]\\
&+\frac1n\sum_i\left[\widehat\mu_{(1)}(X_i)\left(1-\frac{W_i}{\widehat e(X_i)}\right)
-\widehat\mu_{(0)}(X_i)\left(1-\frac{1-W_i}{1-\widehat e(X_i)}\right)\right].
\end{align*}
\inkblue{第一部分是 IPW。}\inkblue{第二部分均值近似为零，因为}
$E[1-W_i/\widehat e(X_i)\mid X_i]\approx0$；条件化后把拟合的结果回归视为固定。
\subsection{半参数效率}
笔记称 AIPW 在非参数估计量中渐近最优。当辅助函数估计足够准确时，
\[
\sqrt n(\widehat\tau_{\rm AIPW}-\widehat\tau^*_{\rm AIPW})\longrightarrow_p0,
\]
星号表示理想 AIPW。该理想估计量只是独立同分布变量的均值，因此其中心极限定理可以传递给使用估计函数的版本：
\notearrow{blue}{$\sqrt n(\widehat\tau_{\rm AIPW}-\widehat\tau^*_{\rm AIPW})\to_p0$（第 12 页）}{理想估计量的独立同分布中心极限定理可用于估计得到的 AIPW。}

\[
\sqrt n(\widehat\tau_{\rm AIPW}-\tau)\Rightarrow N(0,V^*),
\]
\[
V^*=\operatorname{Var}[\tau(X_i)]
+E\left[\frac{b_0^2(X_i)}{1-e(X_i)}\right]
+E\left[\frac{b_1^2(X_i)}{e(X_i)}\right].
\]
\inkblue{页边将其描述为类似 Cramér--Rao 的约束：这是能达到的最好效果／方差。}
\subsection{AIPW 与交叉拟合}
\sourcepages{Lecture1-10.pdf，第 13--14 页。}
将数据随机分成两半 $I_1$ 和 $I_2$。\inkblue{批注：把数据分开，随机选择。}令
\[
\widehat\tau_{\rm AIPW}=\frac{|I_1|}{n}\widehat\tau^{I_1}
+\frac{|I_2|}{n}\widehat\tau^{I_2},
\]
\begin{align*}
\widehat\tau^{I_1}=\frac1{|I_1|}\sum_{i\in I_1}\bigg[&\widehat\mu_{(1)}^{I_2}(X_i)-\widehat\mu_{(0)}^{I_2}(X_i)
+W_i\frac{Y_i-\widehat\mu_{(1)}^{I_2}(X_i)}{\widehat e^{I_2}(X_i)}\\
&-(1-W_i)\frac{Y_i-\widehat\mu_{(0)}^{I_2}(X_i)}{1-\widehat e^{I_2}(X_i)}\bigg].
\end{align*}
渐近分析可以不依赖所选机器学习调整方法的具体形式，而依赖以下条件：
\begin{enumerate}
\item 重叠性：对所有 $x\in\mathcal X$，有 $\eta<e(x)<1-\eta$，其中 $\eta>0$。
\item 辅助函数估计的一致性：
\[
\sup_x|\widehat\mu_{(w)}^{I_2}(x)-\mu_{(w)}(x)|\longrightarrow_p0,
\quad \sup_x|\widehat e^{I_2}(x)-e(x)|\longrightarrow_p0.
\]
\inkblue{蓝字写的是：机器学习调整方法具有大样本一致性。}
\item 结果模型与倾向得分模型的风险衰减：
\[
E[(\widehat\mu_{(w)}^{I_2}(X_i)-\mu_{(w)}(X_i))^2]\,
E[(\widehat e^{I_2}(X_i)-e(X_i))^2]=o(1/n).
\]
\inkblue{若两者均以 $\sqrt n$ 速度一致收敛，则乘积为 $1/n^2$ 阶。一个简单的充分条件是，两种回归调整的均方根误差（RMSE）都以 $o(n^{-1/4})$ 收敛。}
\end{enumerate}
\paragraph{续页上的证明。}
把理想估计量写作
$\widehat\tau^*=|I_1|\widehat\tau^{I_1,*}/n+|I_2|\widehat\tau^{I_2,*}/n$.
要证明 $\sqrt n(\widehat\tau_{\rm AIPW}-\widehat\tau^*)\to_p0$，只需对处理状态下的均值证明相应结果；其余部分类似。定义
\[
\widehat\tau^{I_1}=\widehat\mu^{I_1}_{(1)}-\widehat\mu^{I_1}_{(0)},\qquad
\widehat\mu^{I_1}_{(1)}=\frac1{|I_1|}\sum_{i\in I_1}
\left[\widehat\mu^{I_2}_{(1)}(X_i)+W_i\frac{Y_i-\widehat\mu^{I_2}_{(1)}(X_i)}{\widehat e^{I_2}(X_i)}\right].
\]
减去对应的理想估计量，得到
\begin{align*}
\widehat\mu^{I_1}_{(1)}-\widehat\mu^{I_1,*}_{(1)}
={}&\frac1{|I_1|}\sum_{i\in I_1}
(\widehat\mu^{I_2}_{(1)}(X_i)-\mu_{(1)}(X_i))\left(1-\frac{W_i}{e(X_i)}\right)\\
&+\frac1{|I_1|}\sum_{i\in I_1}W_i(Y_i-\mu_{(1)}(X_i))
\left(\frac1{\widehat e^{I_2}(X_i)}-\frac1{e(X_i)}\right)\\
&-\frac1{|I_1|}\sum_{i\in I_1}W_i(\widehat\mu^{I_2}_{(1)}(X_i)-\mu_{(1)}(X_i))
\left(\frac1{\widehat e^{I_2}(X_i)}-\frac1{e(X_i)}\right).
\end{align*}
\inkblue{蓝字把第一项的界写成 $o_p(1)/n$，最后一项标为 $o_p(1/\sqrt n)$。}
\begin{reviewnote}{original}{S1-011}{把方差的阶标在了余项本身上}{Lecture1-10.pdf，第 14 页}
记第一个均值为 $T_1$，训练数据为 $\mathcal D_2$。由交叉拟合可得 $E[T_1\mid\mathcal D_2]=0$，并且
\[
E[T_1^2\mid\mathcal D_2]
\le \frac{E[(\widehat\mu_1(X)-\mu_1(X))^2\mid\mathcal D_2]}{\eta|I_1|}
=o_p(n^{-1}).
\]
因此 $T_1=o_p(n^{-1/2})$，一般不能说它是 $o_p(n^{-1})$。较小的阶属于其条件方差／二阶矩。第二个均值类似处理，乘积项则使用前面的乘积风险假设。这里指出原批注很可能遗漏了方差运算，同时保留原来的标注。
\end{reviewnote}


对 $K$ 折交叉拟合，在构造第 $i$ 个观测的辅助估计时，排除其所在的第 $k(i)$ 折：
\begin{align*}
\widehat\tau_{\rm AIPW}=\frac1n\sum_i\bigg[&\widehat\mu_{(1)}^{(-k(i))}(X_i)-\widehat\mu_{(0)}^{(-k(i))}(X_i)
+W_i\frac{Y_i-\widehat\mu_{(1)}^{(-k(i))}(X_i)}{\widehat e^{(-k(i))}(X_i)}\\
&-(1-W_i)\frac{Y_i-\widehat\mu_{(0)}^{(-k(i))}(X_i)}{1-\widehat e^{(-k(i))}(X_i)}\bigg].
\end{align*}
\subsection{置信区间}
\sourcepages{Lecture1-10.pdf，第 16、15 页；第 15 页补充第 16 页所用的定义。}
\paragraph{补充：术语。}
原笔记将对数似然的导数标为有效得分：
\[
S(\theta)=\frac{d\log L(\theta)}{d\theta},
\]
其中 $L(\theta)$ 为似然函数。$n$ 个样本得分的经验方差为
\[
\frac1{n-1}\sum_i(S(\theta_i)-\overline{S(\theta)})^2.
\]
有效方差是给定问题下能够达到的最小方差，可以理解为 Cramér--Rao 下界。
\begin{reviewnote}{original}{S1-012}{似然得分与有效影响函数不同}{Lecture1-10.pdf，第 15、16 页}
所写导数是普通的似然得分。有辅助参数时，有效得分还需要去掉辅助得分方向的分量。在 ATE 估计中，下面中心化的 AIPW 单项 $\psi_i^*-\tau$ 是有效影响函数；其方差就是效率界。它不等于原始似然得分，后者的方差是 Fisher 信息。对观测 $O_i$，经验得分方差应使用同一参数值 $\theta$ 下的 $S_\theta(O_i)$，而非不同的参数值 $\theta_i$。第 16 页 AIPW 的经验方差公式，在相应一致性和矩条件下是正确的。
\end{reviewnote}


\paragraph{方差估计。}
定义理想估计量中的单项
\[
\psi_i^*=\mu_{(1)}(X_i)-\mu_{(0)}(X_i)
+\frac{W_i(Y_i-\mu_{(1)}(X_i))}{e(X_i)}
-\frac{(1-W_i)(Y_i-\mu_{(0)}(X_i))}{1-e(X_i)}.
\]
它的经验方差收敛到有效方差：
\[
\frac1{n-1}\sum_i(\psi_i^*-\widehat\tau^*)^2\longrightarrow_p V^*.
\]
页边把 $\widehat\tau^*$ 称为未做交叉拟合的理想 AIPW 估计量。
\notearrow{blue}{未交叉拟合的理想 AIPW 估计（第 16 页）}{$\sum_i(\psi_i^*-\widehat\tau^*)^2/(n-1)$ 中的 $\widehat\tau^*$。}
将每个辅助函数替换为排除第 $k(i)$ 折后训练的估计，得到交叉拟合单项 $\widehat\psi_i$，以及
\[
\widehat V_{\rm AIPW}=\frac1{n-1}\sum_i(\widehat\psi_i-\widehat\tau_{\rm AIPW})^2\longrightarrow_p V^*.
\]
正态分布中央概率区域下画出的置信区间为
\begin{center}\begin{tikzpicture}[x=1.1cm,y=1cm,>=Stealth,font=\small]
\fill[black!12] (-1.4,0)--plot[domain=-1.4:1.4,samples=40](\x,{exp(-\x*\x/2)})--(1.4,0)--cycle;
\draw[->](-2.5,0)--(2.5,0);\draw[domain=-2.5:2.5,samples=60]plot(\x,{exp(-\x*\x/2)});
\draw(-1.4,0)--(-1.4,.4);\draw(1.4,0)--(1.4,.4);\node at(0,.45){$1-\alpha$};
\end{tikzpicture}\end{center}

\begin{reviewnote}{qualification}{S1-013}{置信水平的命名惯例}{Lecture1-10.pdf，第 16 页}
下面公式的渐近覆盖率是 $1-\alpha$；$\alpha$ 为错误概率。手写的“$\alpha$ 水平”可以理解为按显著性水平命名，因此无需改动公式。
\end{reviewnote}

\[
\widehat\tau_{\rm AIPW}\ \pm\ \frac1{\sqrt n}\Phi^{-1}(1-\alpha/2)\sqrt{\widehat V_{\rm AIPW}}.
\]

\section{第 4 讲：估计处理效应的异质性}
\sourcepages{Lecture1-10.pdf，第 17--18 页。}
对于观测 $(X_i,Y_i,W_i)$，其中 $X_i\in\mathcal X$，之前关注的是 ATE $E[Y_i(1)-Y_i(0)]$。异质性分析转而考虑条件平均处理效应（CATE）：
\[
\tau(x)=E[Y_i(1)-Y_i(0)\mid X_i=x].
\]
\notearrow{blue}{$X_i$ 某一类中的平均效应（第 17 页）}{$\tau(x)=E[Y_i(1)-Y_i(0)\mid X_i=x]$。}
\subsection{正则化偏差}
在 $\{Y_i(0),Y_i(1)\}\perp W_i\mid X_i$ 下，
\[
\tau(x)=\mu_{(1)}(x)-\mu_{(0)}(x),\qquad
\mu_{(w)}(x)=E[Y_i\mid X_i=x,W_i=w].
\]
一个直接的 CATE 估计量为 $\widehat\tau_T(x)=\widehat\mu_{(1)}(x)-\widehat\mu_{(0)}(x)$。
\paragraph{个人疑虑。}
\inkblue{如果 $\widehat\mu_{(1)}$ 的样本太少，就可能受到很强的正则化，向一条直线收缩。而 $\widehat\mu_{(0)}$ 样本很多时可能拟合得很好；这样，两者之差 $\widehat\tau_T(x)$ 就会因为前者拟合不准而出问题。}
\subsection{半参数建模}
\[
Y_i(w)=f(X_i)+w\tau(X_i)+\varepsilon_i(w),\qquad
P(W_i=1\mid X_i)=e(X_i),\qquad \tau(x)=\psi(x)\beta.
\]
\inkblue{$f$ 是 $X_i$ 的函数，$\psi(x)$ 是 $X$ 的非参数表示；$\beta$ 是处理效应的系数。}在无混杂性下，Robinson 变换为
\notearrow{black}{$Y_i(w)=f(X_i)+w\tau(X_i)+\varepsilon_i(w)$（第 17 页）}{\inkred{Robinson}：减去 $m(X_i)$，并用 $e(X_i)$ 对 $W_i$ 残差化。}

\[
Y_i-m(X_i)=(W_i-e(X_i))\psi(X_i)\beta+\varepsilon_i,
\quad m(x)=E[Y_i\mid X_i=x]=f(x)+e(x)\tau(x),
\]
其中 $\varepsilon_i=\varepsilon_i(W_i)$。
\notearrow{blue}{$\varepsilon_i=\varepsilon_i(W_i)$（第 17 页）}{$Y_i-m(X_i)=(W_i-e(X_i))\psi(X_i)\beta+\varepsilon_i$ 中的误差。}

构造残差对残差的 OLS 回归，令
\[
\widetilde Y_i^*=Y_i-m(X_i),\qquad
\widetilde Z_i^*=(W_i-e(X_i))\psi(X_i).
\]
将 $\widetilde Y_i^*$ 对 $\widetilde Z_i^*$ 回归，得到 $\widehat\beta_R^*$（原笔记使用的系数符号类似 $\zeta_R$）。\inkblue{这些变换只用到已观测数据和辅助函数；页边称之为数据预处理。}笔记给出
\[
\sqrt n(\widehat\beta_R^*-\beta)\Rightarrow N(0,V_R),
\]
\[
V_R=\operatorname{Var}[\widetilde Z_i^*]^{-1}
\operatorname{Var}[\widetilde Z_i^*\widetilde Y_i^*]
\operatorname{Var}[\widetilde Z_i^*]^{-1}.
\]
同方差时，$\operatorname{Var}(\varepsilon_i\mid X_i,W_i)=b^2$ 为常数，$V_R$ 是 $\beta$ 的半参数有效方差。异方差时，笔记说 $V_R$ 仍适用，但不再是有效方差。
\begin{reviewnote}{original}{S1-014}{Robinson 夹心方差应使用回归残差}{Lecture1-10.pdf，第 18 页}
令列向量 $Z=\widetilde Z_i^*$，$Q=E[ZZ']$，$\varepsilon=\widetilde Y_i^*-Z'\beta$。由正规方程得到影响函数 $Q^{-1}Z\varepsilon$，因此
\[
V_R=Q^{-1}E[ZZ'\varepsilon^2]Q^{-1}.
\]
中间项应为 $\operatorname{Var}(Z\varepsilon)$，而非 $\operatorname{Var}(Z\widetilde Y^*)$。后者包含 $ZZ'\beta$ 的随机波动，而估计分母会抵消这部分波动。这里 $E[Z]=0$，所以 $Q=\operatorname{Var}(Z)$。在条件同方差、$Q$ 非奇异及所需矩有限时，修正后的方差为 $b^2Q^{-1}$。原文关于效率的讨论应针对这个修正后的夹心方差。
\end{reviewnote}


\paragraph{可实施的交叉拟合版本。}
前面的回归是理想形式，因为 $m(X_i)$、$e(X_i)$ 未知。拟合非参数回归 $Y\sim X$ 和 $W\sim X$；构造下式时，\inkblue{排除第 $i$ 个观测所在的折}：
\[
\widetilde Y_i=Y_i-\widehat m^{(-k(i))}(X_i),\qquad
\widetilde Z_i=(W_i-\widehat e^{(-k(i))}(X_i))\psi(X_i).
\]
将 $\widetilde Y_i$ 对 $\widetilde Z_i$ 回归得到 $\widehat\beta_R$。若两个辅助函数估计的 RMSE 均为 $o_p(n^{-1/4})$，笔记得到
\[
\sqrt n(\widehat\beta_R^*-\widehat\beta_R)\longrightarrow_p0.
\]
\subsection{用于处理效应异质性的损失函数}
\sourcepages{Lecture1-10.pdf，第 20、19 页；第 19 页补充第 20 页的残差回归与 R 损失。}
前面的估计量假设 $\tau(x)=\psi(x)\beta$。实际关系往往不是线性的，因此使用机器学习。令
\[
\mu_{(w)}(x)=E[Y(w)\mid X=x],\quad
\varepsilon_i(w)=Y_i(w)-[\mu_{(0)}(X_i)+w\tau(X_i)],
\]
由无混杂性得到 $E[\varepsilon_i(w)\mid X_i,W_i]=0$。\inkblue{这里 $\mu_{(0)}$ 起到前面 $f$ 的作用。}仍有
\[
Y_i-m(X_i)=(W_i-e(X_i))\tau(X_i)+\varepsilon_i,
\quad m(x)=\mu_{(0)}(x)+e(x)\tau(x),\quad \varepsilon_i=\varepsilon_i(W_i).
\]
等价地，
\[
\tau(\cdot)=\arg\min_{t}E[((Y_i-m(X_i))-(W_i-e(X_i))t(X_i))^2].
\]
\inkblue{若 $m$ 和 $e$ 已知，最小化该式即可得到 $\tau(\cdot)$。}
\begin{reviewnote}{qualification}{S1-015}{R 损失的识别与括号中的 RCT}{Lecture1-10.pdf，第 20 页}
总体 R 损失相对最优值的差为
\[
L(t)-L(\tau)=E[e(X)(1-e(X))(t(X)-\tau(X))^2].
\]
因此，只有在 $0<e(X)<1$ 的区域，CATE 才能被唯一确定（忽略零测集）。RCT 中 $e$ 由设计已知，但结果均值 $m(x)=E[Y\mid X=x]$ 一般仍未知，需要估计。因此，括号中提到 RCT，并不意味着两个辅助函数都已知。
\end{reviewnote}
理想的惩罚估计量为
\[
\widehat\tau_R^*(\cdot)=\arg\min_t\left\{
\frac1n\sum_i[(Y_i-m(X_i))-(W_i-e(X_i))t(X_i)]^2+\Lambda_n(t)\right\}.
\]
\notearrow{blue}{对 $\tau(\cdot)$ 的复杂度进行正则化（第 20 页）}{$\Lambda_n(\tau)$：线性情形如 lasso／岭回归；页边还列了 CNN。}\inkblue{通常 $m$ 和 $e$ 未知（原笔记在括号中提到 RCT），因此对它们进行交叉拟合：}
\begin{inkblock}{blue}
\begin{align*}
\widehat L_n(t)&=\frac1n\sum_i
[(Y_i-\widehat m^{(-k(i))}(X_i))-(W_i-\widehat e^{(-k(i))}(X_i))t(X_i)]^2,\\
\widehat\tau_R&=\arg\min_t\{\widehat L_n(t)+\Lambda_n(t)\}.
\end{align*}
对于 lasso，$\widehat\tau(x)=x\widehat\beta$，其中
\[
\widehat\beta=\arg\min_\beta\left\{\frac1n\sum_i
[(Y_i-\widehat m^{(-k(i))}(X_i))-(W_i-\widehat e^{(-k(i))}(X_i))X_i\beta]^2
+\lambda\|\beta\|_1\right\}.
\]
\end{inkblock}

\paragraph{个人补充理解。}
处理效应为常数时，
$Y_i(w)=f(X_i)+\tau w+\varepsilon_i$，且 $E[\varepsilon_i\mid X_i,W_i]=0$；残差估计量为
\[
\widehat\tau_R=\frac{\sum_i(Y_i-\widehat m^{(-k(i))}(X_i))(W_i-\widehat e^{(-k(i))}(X_i))}
{\sum_i(W_i-\widehat e^{(-k(i))}(X_i))^2}.
\tag{原笔记 4.5}
\]
前提是辅助部分收敛足够快。即使 $\tau$ 不为常数，该估计量也会收敛到 $\tau(x)$ 的某个非负加权平均（批注还写了处理组均值减去对照组均值的比较）。
\notearrow{blue}{$\overline Y_{(1)}-\overline Y_{(0)}$（第 19 页页边比较）}{前述 $\tau(x)$ 的非负加权平均。}
异质性较小而重叠性存在困难时，可用式 (4.5) 替代。\notearrow{blue}{常系数残差回归及其加权效应解释（第 19 页）}{用于灵活处理效应函数的 R 损失（第 20 页）。}
\begin{reviewnote}{qualification}{S1-016}{常系数残差回归改变了目标量}{Lecture1-10.pdf，第 19 页}
存在异质性时，其总体极限为
\[
\frac{E[e(X)(1-e(X))\tau(X)]}{E[e(X)(1-e(X))]},
\]
这里要求分母为正，辅助估计收敛足够快。该目标更强调重叠较好的区域。只有加入其他条件，例如效应恒定，或 $e(X)(1-e(X))$ 与 $\tau(X)$ 协方差为零时，它才等于总体 ATE。它不能识别处理由协变量完全确定的区域中的效应。这精确表达了作者关于非负权重的直觉。
\end{reviewnote}


\section{第 5 讲：断点回归设计}
\sourcepages{Lecture1-10.pdf，第 22、21 页；第 21 页解释先前的 IPW 方法为什么不适用。}
使用驱动变量 $Z_i\in\mathbb R$ 引入一个阈值：
\[
W_i=\mathbf1_{\{Z_i\geq c\}},\qquad c\text{ 为阈值。}
\]
\inkblue{当 $Z_i\to c$，靠近 $c$ 的观测特征相似，区别在于是否接受干预。}令 $\mu_{(w)}(z)=E[Y_i(w)\mid Z_i=z]$。若两个条件均值连续，则 $z=c$ 处的 CATE 为
\[
\tau_c=\mu_{(1)}(c)-\mu_{(0)}(c)
=\lim_{z\to c^+}E[Y_i\mid Z_i=z]-\lim_{z\to c^-}E[Y_i\mid Z_i=z].
\]
\paragraph{个人补充：普通 IPW 为什么不适用？}
此前假设无混杂性 $\{Y_i(0),Y_i(1)\}\perp W_i\mid Z_i$ 和重叠性
$\eta\leq P(W_i=1\mid Z_i)\leq1-\eta$。在 RDD 中，无混杂性成立，但重叠性不成立，因为给定 $Z_i=z$ 后，处理概率不是零就是一。\inkblue{蓝字说明：无论如何选择非退化的 $\eta$，都不成立。}必须比较 $c$ 附近的观测，借助相似性而非重叠性。
重叠性不成立意味着不能以 $\sqrt n$ 速度估计 $\tau_c$。其极小极大误差以非参数速度收敛，取决于 $\mu_{(w)}(z)$ 的光滑性。例如，$c$ 处二阶导数一致有界时，可以得到 $n^{-2/5}$ 的估计误差速度。
\begin{reviewnote}{qualification}{S1-017}{RDD 识别及收敛速度结论的适用范围}{Lecture1-10.pdf，第 21、22、24 页}
“不能达到 $\sqrt n$”针对这里的非参数光滑函数类，并非针对所有可能的参数限制。$n^{-2/5}$ 速度需要阈值邻域中的二阶导数界、阈值两侧为正且连续的设计密度、独立观测，以及有界条件方差。只有观测能从两侧趋近阈值时，连续性才识别边界差异。通过平衡平方偏差 $h^4$ 与方差 $(nh)^{-1}$ 的计算是正确的；在 $k$ 阶光滑性及适当的局部多项式拟合下，误差速度为 $n^{-k/(2k+1)}$。
\end{reviewnote}

\subsection{局部线性回归估计}
选择小带宽 $h_n\to0$ 和对称权函数 $K$，在局部拟合两个条件均值。下式中 $W_i$ 的系数
\[
\arg\min_{a,\tau,\beta_{(0)},\beta_{(1)}}\sum_i
K\left(\frac{|Z_i-c|}{h_n}\right)
[Y_i-a-\tau W_i-\beta_{(0)}(Z_i-c)_--\beta_{(1)}(Z_i-c)_+]^2
\]
就是 $\widehat\tau_c$。
\notearrow{red}{权重；负侧；正侧（第 22 页）}{$K(|Z_i-c|/h_n)$；$(Z_i-c)_-$；$(Z_i-c)_+$。}
截距 $a$ 和斜率 $\beta_{(w)}$ 是辅助参数；手写标签指出了负侧和正侧部分。常见核函数包括窗口函数 $K(x)=\mathbf1_{\{|x|\leq1\}}$ 和三角核 $K(x)=(1-|x|)_+$。
\subsection{估计量性质与光滑性}
\sourcepages{Lecture1-10.pdf，第 23--24 页。}
通过以下假设量化连续性：$\mu_{(w)}$ 二次可微，且二阶导数一致有界：
\[
\left|\frac{d^2}{dz^2}\mu_{(w)}(z)\right|\leq B,
\qquad z\in\mathbb R,\quad w\in\{0,1\}.
\]
若光滑性更弱，只有 Lipschitz 连续（原笔记解释为有界但不可微），局部没有明确的线性趋势，局部平均可能更好。若光滑性更强，$k\geq3$ 阶导数有界，则更高阶的局部回归可能改善收敛速度。
\begin{reviewnote}{original}{S1-018}{Lipschitz 连续不等于有界但不可微}{Lecture1-10.pdf，第 23 页}
Lipschitz 连续是指存在有限常数 $L$，使 $|f(x)-f(y)|\le L|x-y|$。它既允许可微函数，也允许 $|x|$ 这样有尖点的函数。仅有有界性不能推出 Lipschitz 连续；而定义域无界时，Lipschitz 函数也不一定有界。作者原本关于较弱光滑性的比较仍有意义。
\end{reviewnote}

\inkred{在 $c$ 处做 Taylor 展开：}
\[
\mu_{(w)}(z)=a_{(w)}+\beta_{(w)}(z-c)+\tfrac12\rho_{(w)}(z-c),
\qquad |\rho_{(w)}(u)|\leq Bu^2.
\]
\inkred{最后一项为余项，}且 $\tau_c=a_{(1)}-a_{(0)}$。分别拟合两侧：
\[
(\widehat a_{(1)},\widehat\beta_{(1)})=arg\min_{a,\beta}
\sum_{Z_i\geq c}K\left(\frac{|Z_i-c|}{h_n}\right)[Y_i-a-\beta(Z_i-c)]^2,
\]
对照侧同理；然后 $\widehat\tau_c=\widehat a_{(1)}-\widehat a_{(0)}$。
\paragraph{窗口核的闭式权重。}
在 $c\leq Z_i\leq c+h_n$ 上，原扫描写作
\[
\widehat a_{(1)}=\sum_i\gamma_iY_i,\qquad
\gamma_i=\frac{\widehat E_{(1)}[(Z_i-c)^2]-\widehat E_{(1)}[Z_i-c](Z_i-c)}
{\widehat E_{(1)}[(Z_i-c)^2]-\widehat E_{(1)}[Z_i-c]^2},
\]
其中 $\widehat E_{(1)}[Z_i-c]$ 是带宽内 $Z_i-c$ 的平均值。还记有
\begin{reviewnote}{original}{S1-019}{局部线性权重的分母缺少样本量}{Lecture1-10.pdf，第 24 页}
令 $N_+=|\{i:c\le Z_i\le c+h_n\}|$、$u_i=Z_i-c$，并在该区间内定义 $m_j=N_+^{-1}\sum_i u_i^j$。归一化的截距权重为
\[
\gamma_i=\frac{m_2-m_1u_i}{N_+(m_2-m_1^2)}.
\]
假设 $m_2-m_1^2>0$。则 $\sum_i\gamma_i=1$ 且 $\sum_i\gamma_i u_i=0$，符合所需约束。原式缺少 $N_+$，其权重和会等于 $N_+$。后面的偏差／噪声计算需要使用归一化权重。
\end{reviewnote}

\[
\sum_i\gamma_i=1,\qquad \sum_i\gamma_i(Z_i-c)=0.
\]

原页上的分解为
\[
\widehat a_{(1)}=a_{(1)}+\underbrace{\sum_i\gamma_i\rho_{(1)}(Z_i-c)}_{\text{\inkred{曲率偏差}}}
+\underbrace{\sum_i\gamma_i(Y_i-\mu_{(1)}(Z_i))}_{\text{\inkred{抽样噪声}}}.
\]
对照侧的截距也有相同形式的分解。
\begin{reviewnote}{original}{S1-020}{Taylor 余项的归一化要前后一致}{Lecture1-10.pdf，第 23、24 页}
按前面的定义，曲率项应为 $\tfrac12\sum_i\gamma_i\rho_{(1)}(Z_i-c)$。也可以重新定义 $r_w=\rho_w/2$，满足 $|r_w(u)|\le Bu^2/2$，然后写为 $\sum_i\gamma_i r_w(u_i)$。遗漏的系数影响常数，不影响所述速度。由于局部线性权重可正可负，曲率偏差的界应使用 $\sum_i|\gamma_i|u_i^2$。
\end{reviewnote}
\notearrow{red}{曲率偏差与抽样噪声（第 24 页）}{\inkred{这两项都需要上界。}}曲率偏差的阶为 $Bh_n^2$，其中 $B$ 是二阶导数的界，$h_n$ 是带宽。抽样噪声均值为零；若 $\operatorname{Var}[Y_i\mid Z_i]\leq b^2$，则其方差不超过 $b^2\sum_i\gamma_i^2$。
若 $Z_i$ 在 $c$ 附近具有连续且非零的密度 $f$，
\[
b^2\sum_i\gamma_i^2\approx\frac{4b^2}{|\{i:c\leq Z_i\leq c+h_n\}|}
\approx\frac{4b^2}{f(c)nh_n}.
\]
因此平方偏差为 $h_n^4$ 阶，方差为 $(nh_n)^{-1}$ 阶，从而
\[
\widehat\tau_c=\tau_c+O_p(n^{-2/5}),\qquad h_n\asymp n^{-1/5}.
\]
\inkred{红字指出 $n^{-2/5}$ 是估计误差，并补充：$k$ 阶光滑性下，$h_n\asymp n^{-1/(2k+1)}$。}

\subsection{驱动变量中的噪声}
\sourcepages{Lecture1-10.pdf，第 25 页。}
\notearrow{blue}{可能有噪声（第 25 页）}{驱动变量 $Z_i$。}引入潜在变量 $U_i\sim G$，并给出两个条件：
\begin{enumerate}
\item 含噪声的驱动变量：原扫描写作 $Z_i\mid U_i\sim N(0,\nu^2)$，其中 $\nu>0$。
\item 驱动变量的无混杂性：$\{Y_i(0),Y_i(1)\}\perp Z_i\mid U_i$。
\end{enumerate}
由 $Y_i=Y_i(W_i)$ 以及 $\{Y_i(0),Y_i(1)\}\perp W_i\mid U_i$，
\[
e(u)=P(W_i=1\mid U_i=u)=P(Z_i\geq c\mid U_i=u)
=1-\Phi\left(\frac{c-u}{\nu}\right),
\]
\[
\alpha_{(w)}(u)=E[Y_i(w)\mid U_i=u],\qquad
\tau(u)=E[Y_i(1)-Y_i(0)\mid U_i=u].
\]
批注定义了标准正态分布函数：
\begin{reviewnote}{original}{S1-021}{含噪声的驱动变量应以潜在变量为条件均值}{Lecture1-10.pdf，第 25 页}
后面的倾向得分及卷积公式要求 $Z_i\mid U_i=u\sim N(u,\nu^2)$，等价于 $Z_i-U_i\mid U_i\sim N(0,\nu^2)$。若对每个 $u$，所写条件均值真的都为零，倾向得分就应为 $1-\Phi(c/\nu)$，与 $u$ 无关。因此原文在条件均值中遗漏了 $u$。后面的分解对该潜在位置模型成立；但仅归一化权重还不能消除混杂偏差，诱导的 $h_+$ 和 $h_-$ 还需要适当平衡。
\end{reviewnote}

\inkred{$\Phi(x)=(2\pi)^{-1/2}\int_{-\infty}^x e^{-t^2/2}\,dt$.}

$U_i$ 无法观测，因此不能直接使用其倾向得分。考虑反卷积型估计量：
\notearrow{black}{$U_i$ 未观测：不能直接使用其倾向得分（第 25 页）}{反卷积权重 $\gamma_+(z),\gamma_-(z)$。}

\[
\widehat\tau_\gamma=\frac1n\sum_{Z_i\geq c}\gamma_+(Z_i)Y_i
-\frac1n\sum_{Z_i<c}\gamma_-(Z_i)Y_i.
\]
选择权重，使
\[
E[\mathbf1_{\{Z_i\geq c\}}\gamma_+(Z_i)]
=E[\mathbf1_{\{Z_i<c\}}\gamma_-(Z_i)]=1.
\]
则
\[
E[\widehat\tau_\gamma]
=\underbrace{\int h_+(u)\tau(u)\,dG(u)}_{\text{\inkred{加权处理效应}}}
+\underbrace{\int(h_+(u)-h_-(u))\alpha_{(0)}(u)\,dG(u)}_{\text{\inkred{混杂偏差}}},
\]
其中权函数为
\[
h_+(u)=\int_c^\infty\gamma_+(z)\varphi_\nu(z-u)\,dz,\qquad
h_-(u)=\int_{-\infty}^c\gamma_-(z)\varphi_\nu(z-u)\,dz,
\]
\begin{inkblock}{blue}
\[
\varphi_\nu(z-u)=\frac1{\sqrt{2\pi\nu^2}}\exp\left[-\frac12\left(\frac{z-u}{\nu}\right)^2\right].
\]
\end{inkblock}

\section{第 6 讲：断点回归的有限样本推断}
\sourcepages{Lecture1-10.pdf，第 27、26、29、28 页；第 26 页承接第 27 页，第 28 页解释第 29 页的权重约束。}
有限样本 RDD 允许不连续／离散的驱动变量，或更复杂的阈值规则。局部线性回归的解 (6.2) 可以写成线性估计量：
\[
\widehat\tau_c(\gamma)=\sum_{i=1}^n\gamma_iY_i.\tag{6.3（原笔记）}
\]
权重仅依赖 $Z_i$。
\notearrow{red}{局部线性解 (6.2)（第 27 页）}{线性估计量 (6.3)，下面继续分析其性质。}
若
\[
Y_i(w)=\mu_{(w)}(Z_i)+\varepsilon_i(w),\qquad
\varepsilon_i(w)\mid Z_i\sim N(0,b^2),
\]
则在给定 $Z_1,\ldots,Z_n$ 后，
\[
\widehat\tau_c(\gamma)\sim N(\widehat\tau_c^*(\gamma),b^2\|\gamma\|_2^2),
\qquad \widehat\tau_c^*(\gamma)=\sum_i\gamma_i\mu_{(W_i)}(Z_i).
\]
这里 $W_i=\mathbf1_{\{Z_i\geq c\}}$，且\inkred{$\|\gamma\|_2=(\gamma_1^2+\cdots+\gamma_n^2)^{1/2}$ 是欧几里得范数。}
\subsection{极小极大线性估计}
\inkred{极小极大是指最小化最坏情形的均方误差（MSE）。}同方差下，线性估计量的条件方差可以直接写为 $b^2\|\gamma\|_2^2$。其偏差则依赖未知的条件均值：
\[
\operatorname{Bias}(\widehat\tau_c(\gamma)\mid Z_1,\ldots,Z_n)
=\sum_i\gamma_i\mu_{(W_i)}(Z_i)-[\mu_{(1)}(c)-\mu_{(0)}(c)].
\]
光滑性限制 $|\mu_{(w)}''(z)|\leq B$ 给出偏差的上界
\[
I_B(\gamma)=\sup_{|\mu_{(w)}''|\leq B}
\left\{\sum_i\gamma_i\mu_{(W_i)}(Z_i)-[\mu_{(1)}(c)-\mu_{(0)}(c)]\right\}.
\]
\inkblue{蓝字说明：因为二阶导数约束的函数类对称，取上确界时不需要绝对值。}因此
\[
|\operatorname{Bias}(\widehat\tau_c(\gamma)\mid Z_1,\ldots,Z_n)|\leq I_B(\gamma),
\]
\[
\operatorname{MSE}(\widehat\tau_c(\gamma)\mid Z_1,\ldots,Z_n)
\leq b^2\|\gamma\|_2^2+I_B^2(\gamma).
\]
\inkred{红字强调“MSE = 方差 + 偏差平方”，并指出这个界不依赖未知的 $\mu$。}左页续写，定义
\[
\gamma^B=\arg\min_\gamma\{b^2\|\gamma\|_2^2+I_B^2(\gamma)\},\qquad
\widehat\tau_c(\gamma^B)=\sum_i\gamma_i^B Y_i.
\tag{原笔记 6.8}
\]
\inkblue{蓝字比较：给定偏差上界 $I_B(\gamma)$ 时，这样得到的 MSE 小于局部线性估计。}
\begin{reviewnote}{qualification}{S1-022}{有限样本极小极大保证的含义}{Lecture1-10.pdf，第 26、27、29、30 页}
精确的条件正态计算需要各单位之间的误差独立且为高斯分布（或使用相应的联合高斯协方差），方差已知，且曲率界 $B$ 有效。仅有边际正态性不够。最小化 $b^2\|\gamma\|_2^2+I_B(\gamma)^2$，是在给定函数类上最小化最坏条件 MSE；并不保证对每一对具体的响应函数，实际 MSE 都更小，也不保证严格改善。只要优化可行，偏差调整后的区间在给定设计条件下有效，包括离散设计。若把估计的 $b$ 或 $B$ 代入，还需要额外论证。
\end{reviewnote}

\subsection{推导极小极大线性权重}
分解为
\[
\mu_{(w)}(z)=a_{(w)}+\beta_{(w)}(z-c)+\rho_{(w)}(z),
\]
其中 $\rho_{(w)}(c)=\rho'_{(w)}(c)=0$，二阶导数以 $B$ 为界。截距和斜率没有限制，因此除非满足下列条件，否则 $I_B(\gamma)$ 可以为无穷大：
\[
\sum_i\gamma_iW_i=1,\qquad \sum_i\gamma_i=0,\qquad
\sum_i\gamma_i(Z_i-c)_+=0,\qquad \sum_i\gamma_i(Z_i-c)_-=0.
\]
加上这些约束后，只需控制余项。将式 (6.8) 改写为
\[
(\gamma^B,t)=\arg\min_{\gamma,t}\{b^2\|\gamma\|_2^2+B^2t^2\},
\]
约束为
\[
\sum_i\gamma_iW_i\rho_{(1)}(Z_i)+\sum_i\gamma_i(1-W_i)\rho_{(0)}(Z_i)\leq t
\]
对所有满足 $\rho_{(w)}(c)=\rho'_{(w)}(c)=0$、$|\rho''_{(w)}(z)|\leq1$ 的 $\rho_{(w)}$ 都成立，同时满足
\[
\sum_i\gamma_iW_i=1,\quad\sum_i\gamma_i=0,\quad
\sum_i\gamma_iW_i(Z_i-c)=0,\quad
\sum_i\gamma_i(1-W_i)(Z_i-c)=0.
\tag{原笔记 6.10}
\]
蓝字提示：求解对应的对偶问题。
\notearrow{blue}{原始极小极大优化问题 (6.10)（第 29 页）}{转化为对偶问题求解。}

\paragraph{个人补充：为什么权重之和为零？}
最终估计的是一个差。将求和拆成处理组和对照组两部分：
\[
\widehat\tau_c(\gamma)=\sum_i\gamma_i^{(1)}Y_i+\sum_i\gamma_i^{(2)}Y_i.
\]
处理组权重之和为 $1$，对照组权重之和为 $-1$，所以
\notearrow{red}{$Y=1$ 组；$Y=0$ 组（第 28 页原标注）}{\inkred{权重分别加总为 $1$、$-1$；总和为 $0$。}}

\begin{inkblock}{red}
\[
\sum_i\gamma_i=1+(-1)=0,\quad
\sum_i\gamma_i^{(1)}=\sum_i\gamma_iW_i=1,\quad
\sum_i\gamma_i^{(2)}=\sum_i\gamma_i(1-W_i)=-1.
\]
\end{inkblock}

图中的权重在阈值左侧近处有负的尖峰，右侧近处有正的尖峰；离阈值更远时，权重改变符号，然后趋近零。
\begin{center}\begin{tikzpicture}[x=1.2cm,y=.8cm,>=Stealth,font=\small]
\draw[->](-2.5,0)--(2.5,0)node[right]{$Z-c$};\draw[->](0,-2)--(0,2.2)node[above]{$\gamma$};
\draw[thick](-2,0).. controls(-1.7,.02)and(-1.5,.8)..(-1.1,.35).. controls(-.5,.1)and(-.2,-.5)..(-.05,-1.9);
\draw[thick](.05,1.9).. controls(.2,.5)and(.45,-.7)..(.9,-.45).. controls(1.4,-.3)and(1.6,.0)..(1.9,0);
\end{tikzpicture}\end{center}
原页红字标为 $Y=1$ 和 $Y=0$，下面的公式则使用处理指示变量 $W$。
\begin{reviewnote}{original}{S1-023}{红色分组标签指处理而非结果}{Lecture1-10.pdf，第 28 页}
这些标签应为 $W=1$ 和 $W=0$，结果 $Y$ 不一定是二元变量。权重和的约束本身正确：处理组权重和为 $1$，对照组为 $-1$。
\end{reviewnote}

\subsection{线性估计量的推断}
\sourcepages{Lecture1-10.pdf，第 30--31 页。}
构造 $\widehat\tau_c(\gamma^B)$ 的置信区间。令
\[
\mathrm{err}=\widehat\tau_c-\tau_c,\qquad
\mathrm{err}\mid Z_1,\ldots,Z_n\sim N(\mathrm{bias},b^2\|\gamma^B\|_2^2),
\]
其中由式 (6.10) 得到 $|\mathrm{bias}|\leq Bt$。根据高斯分布的单峰性，对于 $S\sim N(0,1)$，
\[
P(|\mathrm{err}|\geq\zeta)\leq P(|Bt+b\|\gamma^B\|_2S|\geq\zeta).
\]
条件覆盖率至少为 $1-\alpha$ 的置信区间是
\[
I_\alpha=(\widehat\tau_c(\gamma^B)-\zeta_\alpha^B,
\widehat\tau_c(\gamma^B)+\zeta_\alpha^B),
\]
\[
\zeta_\alpha^B=\inf\{\zeta:P(|Bt+b\|\gamma^B\|_2S|>\zeta)\leq\alpha\}.
\tag{原笔记 6.13}
\]
\inkblue{蓝字强调：条件于 $Z_i$，因此驱动变量服从任何分布都可以，包括非标准的 RDD 情形。}
\paragraph{离散驱动变量。}
单独限制 $|\mu_{(w)}''(z)|\leq B$ 不一定能识别 $\tau_c$，但式 (6.13) 仍有效。$Z_i$ 连续时，区间以 $n^{-2/5}$ 速度缩小；$Z_i$ 离散时，长度未必趋于零。\inkblue{当 $c$ 不在驱动变量的支持集内，需要用满足二阶导数界的插值函数表示 $\mu_{(w)}(z)$。}
\paragraph{多维驱动变量。}
例如二维、三维区域以及地理边界。不能直接使用一维阈值 RDD；\inkblue{有的方法将其化为一维距离，笔记认为这样不太稳定。}极小极大方法则可以直接推广到多维。
\subsection{超出同方差情形}
当误差 $\varepsilon_i=Y_i-\mu_{(W_i)}(Z_i)$ 不再是方差为 $b^2$ 的同方差高斯误差时，式 (6.11) 的精确正态分布结论不成立。利用中心极限定理来说明
\[
\widehat\tau_c(\gamma)\mid Z_1,\ldots,Z_n
\ \approx\ N\left(\widehat\tau_c^*(\gamma),
\sum_i\gamma_i^2\operatorname{Var}[Y_i\mid Z_i,W_i]\right).
\tag{原笔记 6.14}
\]
若高斯近似有效，仍可构造置信区间。将条件方差估计为
\[
\widehat V_n=\sum_i\gamma_i^2(Y_i-\widehat\mu_{(W_i)}(Z_i))^2,
\]
其中拟合均值来自局部线性回归。
\notearrow{blue}{局部线性拟合的均值（第 31 页）}{残差方差估计量中的 $\widehat\mu_{(W_i)}(Z_i)$。}
\notearrow{blue}{如果均值拟合错了（第 31 页）}{\inkblue{残差会膨胀，结果过于保守。}}极小极大准则 (6.8) 针对同方差；异方差下，将各个方差 $b_i^2$ 引入用于构造对偶问题的式 (6.10)。
\begin{reviewnote}{qualification}{S1-024}{残差方差估计不会自动保守}{Lecture1-10.pdf，第 31 页}
对于独立的求值观测和固定的拟合均值，残差平方的期望等于条件方差加均值误差平方。但若回归在同一样本内拟合，过拟合和杠杆作用反而可能使残差缩小。因此，保守性需要适当的留一法／交叉拟合方差估计，或控制样本内拟合的理论。条件高斯近似还需要加权 Lindeberg 条件（没有占主导地位的观测），方差估计需要一致或保守有效。这些要求不能仅靠把 $b^2$ 换成 $b_i^2$ 来保证。
\end{reviewnote}

\section{第 7 讲：平衡估计量}
\sourcepages{Lecture1-10.pdf，第 32--33 页。}
标题将平衡性与估计 $\widehat e(x)$ 的原理联系起来。
\notearrow{black}{平衡性／$\widehat e(x)$ 的设计原理（第 32 页）}{改善倾向得分与回归函数；高维 ATE 估计。}
记录的两个目标是：改善倾向得分和回归函数，以及在高维协变量下估计 ATE。IPW、AIPW 都使用倾向得分；这一讲通过凸优化推导如何估计 $e(x)$。
\subsection{回顾 IPW}
给定独立同分布观测 $(X_i,Y_i,W_i)\in\mathcal X\times\mathbb R\times\{0,1\}$，目标为 $\tau=\mu_{(1)}-\mu_{(0)}$，其中 $\mu_{(w)}=E[Y_i(w)]$。假设
\[
\{Y_i(0),Y_i(1)\}\perp W_i\mid X_i,\qquad
0<\eta\leq e(X_i)\leq1-\eta<1,
\]
其中 $e(x)=P(W_i=1\mid X_i=x)$，$\eta$ 为正常数。用下式估计处理状态下的均值：
\[
\widehat\mu_{\rm IPW}(1)=\frac1n\sum_i\frac{W_iY_i}{\widehat e(X_i)}.
\tag{原笔记 7.3}
\]
笔记将下面的一致性论证标为\inkblue{“简略，不重要”}，但仍记录了三步：
\begin{enumerate}
\item 总体平衡：无混杂性下，$\mu_{(1)}=E[W_iY_i/e(X_i)]$。
\item 理想估计量：使用真实的 $e(x)$ 时无偏；重叠性给出有限方差。
\item 可实施近似：若估计的倾向得分也满足重叠性，由 Cauchy--Schwarz 不等式可得
\begin{align*}
|\widehat\mu_{\rm IPW}(1)-\widehat\mu^*_{\rm IPW}(1)|
&\leq\left[\frac1n\sum_i\left(\frac1{\widehat e(X_i)}-\frac1{e(X_i)}\right)^2\right]^{1/2}
\left[\frac1n\sum_i(W_iY_i)^2\right]^{1/2}\\
&\leq\frac1{\eta^2}\left[\frac1n\sum_i(\widehat e(X_i)-e(X_i))^2\right]^{1/2}
\left[\frac1n\sum_i(W_iY_i)^2\right]^{1/2}.
\end{align*}
\inkblue{页边回顾 $|a\cdot b|^2\leq\|a\|^2\|b\|^2$，并补充：$n\to\infty$ 时，相应的 AIPW 误差相对于 $1/\sqrt n$ 可以忽略。}
\end{enumerate}
\subsection{总体平衡与样本平衡}
用预先指定的基展开条件均值：
\[
\mu_{(w)}(x)=\sum_{j=1}^\infty\beta_{j(w)}\psi_j(x).
\]
\notearrow{red}{多项式基或 Fourier 基（第 33 页）}{条件均值展开中的 $\psi_j(x)$。}在合理的正则条件下，
\[
\mu_{(w)}=\sum_{j=1}^\infty\beta_{j(w)}E[\psi_j(X_i)].
\]
\notearrow{black}{两边取期望（第 33 页）}{$\mu_{(w)}=\sum_j\beta_{j(w)}E[\psi_j(X)]$。}
\inkred{这些正则条件会不会有问题？}
\begin{reviewnote}{qualification}{S1-025}{回应基展开的正则性问题}{Lecture1-10.pdf，第 33 页}
这个问题是合理的：级数逐点收敛本身不足以交换求和与期望。\par
一个充分条件是 $\sum_j|\beta_{j(w)}|E|\psi_j(X)|<\infty$；或者要求部分和在 $L^1$ 中收敛。概率分布下的 $L^2$ 收敛也能推出 $L^1$ 收敛。满足重叠性时，同样的论证可以控制加权期望。总体恒等式也可以直接通过条件于 $X$ 来证明，不必使用无限基展开。
\end{reviewnote}

在无混杂性下，$Y_i=\mu_{(W_i)}(X_i)+\varepsilon_i$，且 $E[\varepsilon_i\mid X_i,W_i]=0$。总体 IPW 论证为
\begin{align*}
E\left[\frac{W_iY_i}{e(X_i)}\right]
&=E\left[\frac{W_i}{e(X_i)}\sum_j\beta_{j(1)}\psi_j(X_i)\right]\\
&=\sum_j\beta_{j(1)}E\left[\frac{W_i\psi_j(X_i)}{e(X_i)}\right]
=\sum_j\beta_{j(1)}E[\psi_j(X_i)]=\mu_{(1)}.
\end{align*}
\notearrow{red}{核心：$E[W_i\mid X_i]=e(X_i)$（第 33 页）}{$E[W_i\psi_j(X_i)/e(X_i)]=E[\psi_j(X_i)]$。}因此选择权重，使样本中满足
\[
\frac1n\sum_i\frac{W_i\psi_j(X_i)}{\widehat e(X_i)}
\approx\frac1n\sum_i\psi_j(X_i),\qquad j=1,2,\ldots.
\tag{原笔记 7.8}
\]
\subsection{倾向得分估计的平衡损失}
\sourcepages{Lecture1-10.pdf，第 35、34 页；第 34 页补充高维界中用到的范数知识。}
假设结果模型为线性：$\mu_{(w)}(x)=x\beta_{(w)}$，倾向得分模型为 logistic：
\[
e(x)=\frac1{1+\exp(-x\theta)}.
\]
\inkblue{蓝字说明：此时样本平衡只涉及协变量 $X_i$，平衡的目标是均值。}式 (7.8) 变为
\[
\frac1n\sum_i(1+e^{-X_i\widehat\theta})W_iX_i=\frac1n\sum_iX_i.
\tag{原笔记 7.9}
\]
这是以下凸优化问题的 KKT 条件：
\[
\widehat\theta=\arg\min_\theta\frac1n\sum_i\ell_\theta(X_i,Y_i,W_i),\qquad
\ell_\theta(X_i,Y_i,W_i)=W_ie^{-X_i\theta}+(1-W_i)X_i\theta.
\tag{原笔记 7.10}
\]
因此可以用 Newton 下降法求解。
\notearrow{black}{平衡方程 (7.9) 是凸损失的 KKT 条件（第 35 页）}{Newton 下降法；原笔记称解唯一。}
手写文字称其有唯一解，并用它作为损失函数。
\begin{reviewnote}{original}{S1-026}{凸性本身不能保证存在唯一有限解}{Lecture1-10.pdf，第 35 页}
Hessian 为 $n^{-1}\sum_iW_i e^{-X_i\theta}X_i'X_i$。它是半正定的；只有处理组协变量满足适当的张成／秩条件时才正定。有限极小值点还需要存在性；若目标协变量均值位于处理组可达到的凸支持集之外，精确平衡可能无法实现。因此，仅有凸性既不能保证唯一性，也不能保证极小值能取到，Newton 迭代还需要通常的保护措施。所写一阶平衡方程是正确的。
\end{reviewnote}
\inkblue{页边写道：同时估计 $\mu_{(0)}$ 和 $\mu_{(1)}$，需要两个倾向得分模型。}
\begin{reviewnote}{qualification}{S1-027}{两次校准拟合与一个真实倾向得分}{Lecture1-10.pdf，第 35 页}
处理机制只有一个倾向得分 $e(x)$，对照概率为 $1-e(x)$。在这一构造中，分别做两次校准拟合，是实现各处理组平衡的一种方式；但估计两个均值并不在数学上要求存在两个不同的真实倾向得分模型。
\end{reviewnote}

\subsection{高维混杂变量下的 ATE 估计}
设 $X_i\in\mathbb R^p$，可能有 $p\gg n$，且 $\mu_{(w)}(x)=x\beta_{(w)}$。高维下很难充分准确地估计 $e(x)$，精确权重无法得到，因此寻求近似平衡。由 Hölder 不等式，
\begin{align*}
&\frac1n\sum_i\frac{W_i\mu_{(1)}(X_i)}{\widehat e(X_i)}-\frac1n\sum_i\mu_{(1)}(X_i)\\
&\quad=\left(\frac1n\sum_i\frac{W_iX_i}{\widehat e(X_i)}-\frac1n\sum_iX_i\right)\beta_{(1)}\\
&\quad\leq\left\|\frac1n\sum_i\frac{W_iX_i}{\widehat e(X_i)}-\frac1n\sum_iX_i\right\|_\infty\|\beta_{(1)}\|_1.
\end{align*}
\inkred{第一个范数标为不平衡程度。}\inkblue{手写补充回顾了 Hölder 不等式：}
\inkblue{$\sum_i|a_ib_i|\leq(\sum_i|a_i|^p)^{1/p}(\sum_i|b_i|^q)^{1/q}$.}
\paragraph{个人理解与问题。}
\inkblue{好的逆概率权重使最坏情形下的不平衡程度尽可能小。}还剩两个问题：（1）如何参数化 $\widehat e(X_i)$，实现上确界范数下的近似平衡；（2）只有 $\|\beta_{(1)}\|_1$ 有界时，该界才有意义。\inkblue{高维统计通常假设系数向量稀疏，而不是直接限制其 $1$ 范数。}
\paragraph{补充：范数回顾。}
左页将“$0$ 范数”列为 $\sum_iX_i^0$，$1$ 范数为 $\sum_i|X_i|$，$2$ 范数为 $(\sum_iX_i^2)^{1/2}$，无穷范数为 $\max_i|X_i|$。一般表达式写为 $\|x\|_p=(\sum_i x_i^p)^{1/p}$。
\begin{reviewnote}{original}{S1-028}{修正范数简写并明确 Hölder 共轭指数}{Lecture1-10.pdf，第 34、35 页}
应使用 $\|x\|_0=\sum_i\mathbf1\{x_i\ne0\}$（这是计数，不是范数），以及 $p\ge1$ 时的 $\|x\|_p=(\sum_i|x_i|^p)^{1/p}$。$x_i^0$ 不能正确处理零分量；对负分量取奇数或非整数次幂，也不能定义所需的实值范数。Hölder 不等式要求 $p,q\ge1$ 且 $1/p+1/q=1$，包括端点 $(1,\infty)$。所写不平衡界由该端点不等式得到；加绝对值后给出的是误差大小的上界。
\end{reviewnote}


\subsection{近似平衡权重与回归增广}
\sourcepages{Lecture1-10.pdf，第 36--37 页。}
选择合适的调节参数 $\zeta$，求解
\[
\widehat\gamma=\arg\min_{\gamma_i\geq1}
\left\{\frac1{n^2}\|\gamma\|_2^2+
\zeta\left\|\frac1n\sum_i(\gamma_iW_i-1)X_i\right\|_\infty^2\right\}.
\tag{原笔记 7.12}
\]
\notearrow{red}{$L_2$ 大小的平方；不平衡程度；调节权重（第 36 页）}{$n^{-2}\|\gamma\|_2^2$；无穷范数的平方；$\zeta$。}然后将 $1/\widehat e(X_i)=\widehat\gamma_i$ 作为逆倾向得分权重。利用次高斯集中不等式，可对理想权重验证如下形式的界：
\[
\frac1{n^2}\|1/e(X_i)\|_2^2\leq\frac{\eta^{-2}}n,
\qquad
\left\|\frac1n\sum_i(\gamma_iW_i-1)X_i\right\|_\infty^2
=O_p\left(\frac{\eta^{-2}\log p}{n}\right).
\]
笔记说 $\widehat e(X_i)$ 不需要很好地估计 $e(X_i)$；选择合适的 $\zeta$ 后，这些逆权重仍满足偏差界。
接着用高维回归增广加权估计量，以去掉对 $\|\beta_{(1)}\|_1$ 的依赖。增广平衡估计量为
\[
\widehat\mu_{\rm AB}(1)=\frac1n\sum_i
\left[\widehat\mu_{(1)}(X_i)+\widehat\gamma_iW_i(Y_i-\widehat\mu_{(1)}(X_i))\right],
\quad\widehat\mu_{(1)}(X_i)=X_i\widehat\beta_{(1)}.
\tag{原笔记 7.14}
\]
其分解为
\begin{align*}
\widehat\mu_{\rm AB}(1)={}&\underbrace{\frac1n\sum_iX_i\beta_{(1)}}_{\textcolor{red}{\mu_{(1)}(X_i)}\text{\inkred{ 的均值}}}
+\underbrace{\frac1n\sum_i\widehat\gamma_iW_i(Y_i-X_i\beta_{(1)})}_{\text{\inkred{零均值噪声}}}\\
&+\left[\frac1n\sum_i(1-\widehat\gamma_iW_i)X_i\right](\widehat\beta_{(1)}-\beta_{(1)}).
\end{align*}
最后一项的绝对值上界为
\begin{inkblock}{red}
\[
\left\|\frac1n\sum_i(1-\widehat\gamma_iW_i)X_i\right\|_\infty
\|\widehat\beta_{(1)}-\beta_{(1)}\|_1.
\]
\end{inkblock}
\notearrow{red}{权重只需修正 $\beta_{(1)}-\widehat\beta_{(1)}$，不必修正整个系数向量（第 36 页）}{不平衡程度乘回归误差的余项。}高维回归给出该误差的界：在稀疏性 $\|\beta_{(1)}\|_0\leq k$ 与限制特征值条件下，lasso 可达到
\[
\|\widehat\beta_{(1)}-\beta_{(1)}\|_1
=O_p\left(k\sqrt{\frac{\log p}{n}}\right).
\]
\paragraph{将步骤组合起来。}
假设设计满足上述界所需的条件，$\beta_{(1)}$ 为 $k$ 稀疏，且 $k\ll\sqrt n/\log p$。
\begin{enumerate}
\item 在处理组上拟合 lasso；系数的 $1$ 范数误差为 $o_p(1/\sqrt{\log p})$。
\item 用式 (7.12) 拟合 $\widehat\gamma$，得到无穷范数不平衡程度 $O_p(\sqrt{\log(p)/n})$。
\item 使用式 (7.14)，将这些界代入分解：
\[
\widehat\mu_{\rm AB}(1)-\frac1n\sum_iX_i\beta_{(1)}
=\frac1n\sum_i\widehat\gamma_iW_i(Y_i-X_i\beta_{(1)})+o_p(n^{-1/2}).
\tag{原笔记 7.17}
\]
\inkred{噪声项为 $O_p(n^{-1/2})$。}
\begin{reviewnote}{qualification}{S1-029}{高维平衡推断所需的条件}{Lecture1-10.pdf，第 32、36、37 页}
重叠性本身不能保证有限方差，还需 $E[Y(w)^2]<\infty$。集中界也使用了受控的次高斯协变量，而 lasso 速度需要合适的调节参数、限制特征值以及噪声假设。将两个速度相乘可见，所写稀疏性条件使余项为 $o_p(n^{-1/2})$。条件于 $(X,W)$ 后，权重必须只依赖这些变量，而不依赖结果噪声。中心极限定理还要求不存在占主导地位的加权残差；所写学生化分母使用了未知真实值 $\beta_1$，因此可实施推断还需论证用估计值替换它的有效性。
\end{reviewnote}

\end{enumerate}
还成立中心极限定理：
\[
\frac{\sum_i\widehat\gamma_iW_i(Y_i-X_i\beta_{(1)})}
{\sqrt{\sum_i\widehat\gamma_i^2W_i^2(Y_i-X_i\beta_{(1)})^2}}
\Rightarrow N(0,1).
\]
\paragraph{个人保留意见。}
\begin{inkblock}{blue}
这比预想的弱：只给出了样本均值
$\overline\mu_{(1)}=n^{-1}\sum_iX_i\beta_{(1)}$ 的置信区间，而不是 $\mu_{(1)}$ 本身的区间。要得到总体均值的结果，需要证明估计权重 $\widehat\gamma_i$ 的收敛性。
\end{inkblock}

\begin{reviewnote}{qualification}{S1-030}{样本目标与总体目标的不确定性}{Lecture1-10.pdf，第 37 页}
针对总体目标，还要计入抽样波动 $n^{-1}\sum_i\mu_1(X_i)-E[\mu_1(X)]$。拟合权重收敛到稳定函数，是建立影响函数结论的一条有用的充分路径，但并非所有构造中逻辑上必需的条件。必须建立联合极限及其方差。作者指出所写条件区间只覆盖样本均值，这一提醒是正确的。
\end{reviewnote}


\section{第 8 讲：面板数据方法}
\sourcepages{Lecture1-10.pdf，第 39、38 页；第 38 页补充面板布局和单位内部的方差。}
面板（纵向）数据在时刻 $t=1,\ldots,T$ 观测单位 $i=1,\ldots,n$。左页示意图为矩阵：各行对应单位，各列对应时间点 $t_1,t_2,\ldots,T$。
\[
\begin{array}{c|cccc}
 &t=t_1&t=t_2&\cdots&t=T\\\hline
 i=1&&&&\\i=2&&&&\\\vdots&&&&\\i=n&&&&
\end{array}
\]

\subsection{常数处理效应模型}
结果为 $Y_{it}\in\mathbb R$，处理分配为 $W_{it}\in\{0,1\}$。基本模型是
\[
Y_{it}=Y_{it}(0)+W_{it}\tau,\qquad i=1,\ldots,n,\quad t=1,\ldots,T.
\tag{原笔记 8.1}
\]
这个非常基础的模型没有考虑处理效应异质性及动态效应。笔记解释：所有单位、所有时点的处理效应相同；单位在 $t$ 时点的结果只依赖该时点的处理，不依赖之前的处理。\inkred{$\tau$ 不依赖 $i,t$；过去的时点不进入模型。}
\subsection{双向模型}
\[
Y_{it}=\alpha_i+\beta_t+W_{it}\tau+\varepsilon_{it},
\qquad E[\varepsilon\mid\alpha,\beta,W]=0.
\tag{原笔记 8.2}
\]
\inkred{$\alpha_i$、$\beta_t$ 分别为单位和时间的平移。}\inkred{红字将该条件称为处理的严格外生性，比无混杂性更强。}
\begin{reviewnote}{qualification}{S1-031}{严格外生性与无混杂性的范围不同}{Lecture1-10.pdf，第 39 页}
这里严格外生性要求：给定整个处理路径和固定效应后，误差均值为零；处理决策不能响应任何时期的个体特异误差。它比相应的同期均值外生性更强，但不能笼统地与所有分布意义上的无混杂性假设排列强弱。应明确条件集合，以及讨论的是条件均值还是完整的独立性。
\end{reviewnote}
例如有两个时期，一组的处理路径为 $(0,0)$，另一组为 $(0,1)$。双重差分的闭式估计为
\notearrow{blue}{DID（第 39 页）}{两组各自的前后差，再做组间差。}

\[
\widehat\tau=\frac1{|\{i:W_{i2}=1\}|}\sum_{W_{i2}=1}(Y_{i2}-Y_{i1})
-\frac1{|\{i:W_{i2}=0\}|}\sum_{W_{i2}=0}(Y_{i2}-Y_{i1}).
\]
\paragraph{笔记中的评论。}
局限是：当 $n,T$ 很大且处理分配分布一般时，模型 (8.2) 可能不成立。类似地，多于两个时期的 OLS 有时也称为 DID，但原笔记认为 OLS 没有闭式解。列出的优点是平行趋势要求比较直观（单位只在 $\alpha_i$ 上存在差异）；没有平行趋势就不能使用。
\begin{reviewnote}{original}{S1-032}{一般双向 OLS 仍有闭式解}{Lecture1-10.pdf，第 39 页}
对于平衡面板，将 $Y$ 和 $W$ 对单位及时间虚拟变量残差化：$\widetilde W_{it}=W_{it}-\overline W_{i\cdot}-\overline W_{\cdot t}+\overline W_{\cdot\cdot}$，$Y$ 同理。当 $\sum_{it}\widetilde W_{it}^2>0$ 时，
\[
\widehat\tau=\frac{\sum_{it}\widetilde W_{it}\widetilde Y_{it}}{\sum_{it}\widetilde W_{it}^2}.
\]
因此，简单的两组前后差表达式未必适用，但 OLS 系数仍有闭式解。固定效应模型意味着条件均值趋势平行，而不是含噪声的实际路径完全平行。
\end{reviewnote}


\paragraph{补充：噪声建模与方差估计。}
将各行向量 $\varepsilon_i=(\varepsilon_{i1},\ldots,\varepsilon_{iT})$ 视为在单位之间独立同分布，并有 $\operatorname{Var}[\varepsilon_i]=\Sigma$。假设每个 $\varepsilon_{it}$ 都相互独立，很难论证，也不合理。$T=2$ 时，直接估计差分的方差：
\[
\widehat{\operatorname{Var}}[\widehat\tau]
=\frac{\widehat{\operatorname{Var}}[Y_{i2}-Y_{i1}\mid W_{i2}=1]}{|\{i:W_{i2}=1\}|}
+\frac{\widehat{\operatorname{Var}}[Y_{i2}-Y_{i1}\mid W_{i2}=0]}{|\{i:W_{i2}=0\}|}.
\tag{原笔记 8.4}
\]
同一行／单位内的误差相关时，该方法仍稳健。
% 后续由 stats1-tail-zh.tex 翻译原扫描第40--58页，从交互面板模型开始。
